\documentclass[a4paper,11pt]{article}

% Packages
\usepackage{amsmath, amssymb, amsthm}
\usepackage{algorithm}
\usepackage{algpseudocode}
\usepackage{enumitem}
\usepackage{geometry}
\geometry{margin=2.5cm}

% Optional: nicer section spacing
\usepackage{titlesec}
\titleformat{\section}{\large\bfseries}{\thesection}{0.5em}{}

\begin{document}

	\begin{center}
		{\Large \textbf{CSE2310 Dynamic Programming Peer Review Task a}}\\[4pt]
		Student number: \textbf{6152937}
	\end{center}

	\vspace{1em}

	\section{Rhoda's luggage problem}

	\subsection{Problem statement}
	We have $n$ souvenirs. Item $i$ has volume $l_i$ (liters), liquid amount $k_i$ (liters), and value $v_i$ (euros).
	The suitcase has capacity $L$ liters. Up to $K$ liters of liquid are free. If total liquid is $q$, then the extra cost is
	\[
	\mathrm{cost}(q) = c \cdot \max\Bigl(0,\Bigl\lceil \tfrac{q}{K}\Bigr\rceil - 1\Bigr).
	\]
	The goal is to maximize net value: total value minus $\mathrm{cost}(\text{total liquid})$. \ \ :contentReference[oaicite:0]{index=0}

	\subsection{1. Small instance with $n=3$ and solution}
	Let $L=6$, $K=3$, $c=5$, and the three items be:
	\[
	\begin{array}{c|ccc}
		i & l_i & k_i & v_i\\ \hline
		1 & 3 & 3 & 12\\
		2 & 3 & 2 & 10\\
		3 & 2 & 0 & 4
	\end{array}
	\]
	Enumerate feasible subsets (volume at most $6$):

	\begin{itemize}[leftmargin=2em]
		\item $\{1\}$: volume $3$, liquid $3$, value $12$, cost $0$, net $12$.
		\item $\{2\}$: volume $3$, liquid $2$, value $10$, cost $0$, net $10$.
		\item $\{3\}$: volume $2$, liquid $0$, value $4$, cost $0$, net $4$.
		\item $\{1,2\}$: volume $6$, liquid $5$, value $22$.
		Here $\lceil 5/3\rceil = 2$, so cost $= (2-1)c = 5$. Net $=22-5=17$.
		\item $\{1,3\}$: volume $5$, liquid $3$, value $16$, cost $0$, net $16$.
		\item $\{2,3\}$: volume $5$, liquid $2$, value $14$, cost $0$, net $14$.
		\item $\{1,2,3\}$: volume $8$ exceeds $L$, infeasible.
	\end{itemize}

	Optimal is $\{1,2\}$ with net value $17$. It is optimal to buy extra liquid allowance (one extra block of $K$ liters).

	\subsection{2. If $k_i=0$ for all items}
	If $k_i=0$ for all $i$, then total liquid is always $0$, so $\mathrm{cost}(0)=0$ and the problem becomes the classic $0$-$1$ knapsack:
	maximize $\sum v_i$ subject to $\sum l_i \le L$. \ \ :contentReference[oaicite:1]{index=1}

	\subsection{3. Example where buying extra liquid is optimal, and where it is not}
	\paragraph{Buying extra liquid is optimal.}
	The instance from Question 1 is such an example: choosing $\{1,2\}$ forces liquid $5>K$, costs $5$, and still beats any solution without buying (net $17$ vs best without buying net $16$). \ \ :contentReference[oaicite:2]{index=2}

	\paragraph{Buying extra liquid is not optimal.}
	Let $L=6$, $K=3$, $c=20$, and:
	\[
	\begin{array}{c|ccc}
		i & l_i & k_i & v_i\\ \hline
		1 & 3 & 3 & 12\\
		2 & 3 & 2 & 10\\
		3 & 3 & 3 & 11
	\end{array}
	\]
	Feasible options:
	\begin{itemize}[leftmargin=2em]
		\item Without buying, best is $\{1\}$ with net $12$ (or $\{3\}$ with net $11$).
		\item If we buy extra liquid: $\{1,2\}$ has liquid $5$, so $\lceil 5/3\rceil=2$ and cost $20$, net $22-20=2$.
		Also $\{1,3\}$ has liquid $6$, $\lceil 6/3\rceil=2$, cost $20$, net $23-20=3$.
	\end{itemize}
	So it is not optimal to buy extra liquid. \ \ :contentReference[oaicite:3]{index=3}

	\subsection{4. Recursive function (math notation) for the maximum net value}
	Define a helper DP that ignores the step cost until the end:

	Let $V(i,\ell,q)$ be the maximum total value achievable using only items $\{1,\dots,i\}$ with
	total volume at most $\ell$ and total liquid exactly $q$.
	If impossible, $V(i,\ell,q)=-\infty$.

	Base:
	\[
	V(0,\ell,0)=0 \ \text{for all } \ell \in \{0,\dots,L\},\qquad
	V(0,\ell,q)=-\infty \ \text{for } q>0.
	\]

	Recurrence for $i\ge 1$:
	\[
	V(i,\ell,q) =
	\max\Bigl(
	V(i-1,\ell,q), \;
	v_i + V(i-1,\ell-l_i,q-k_i)
	\Bigr),
	\]
	where the second term is only allowed if $\ell\ge l_i$ and $q\ge k_i$.

	Given $V$, the maximum net value is:
	\[
	\mathrm{OPT}(n,L)=
	\max_{\ell \in \{0,\dots,L\}}
	\max_{q \in \{0,\dots,Q\}}
	\Bigl(
	V(n,\ell,q) - c\cdot \max(0,\lceil q/K\rceil - 1)
	\Bigr),
	\]
	where $Q=\sum_{i=1}^n k_i$ (and $Q \le nK$ since $k_i \le K$).

	We call $\mathrm{OPT}(n,L)$ to obtain the optimal net value for the full instance. \ \ :contentReference[oaicite:4]{index=4}

	\subsection{5. Efficient iterative DP algorithm (pseudocode)}
	\begin{algorithm}[H]
		\caption{Maximum net value with liquid cost blocks}
		\begin{algorithmic}[1]
			\State \textbf{input:} $n,L,K,c$ and arrays $(l_i,k_i,v_i)$ for $i=1..n$
			\State \textbf{output:} maximum net value
			\State $Q \gets \sum_{i=1}^n k_i$
			\State Create array $dp[0..L][0..Q]$ and fill with $-\infty$
			\State $dp[0][0] \gets 0$
			\For{$i \gets 1$ to $n$}
			\State Create array $ndp \gets dp$ \Comment{not taking item $i$}
			\For{$\ell \gets 0$ to $L$}
			\For{$q \gets 0$ to $Q$}
			\If{$dp[\ell][q] > -\infty$}
			\If{$\ell + l_i \le L$ \textbf{and} $q + k_i \le Q$}
			\State $ndp[\ell + l_i][q + k_i] \gets \max(ndp[\ell + l_i][q + k_i],\ dp[\ell][q] + v_i)$
			\EndIf
			\EndIf
			\EndFor
			\EndFor
			\State $dp \gets ndp$
			\EndFor
			\State $best \gets -\infty$
			\For{$\ell \gets 0$ to $L$}
			\For{$q \gets 0$ to $Q$}
			\If{$dp[\ell][q] > -\infty$}
			\State $blocks \gets \left\lceil \frac{q}{K} \right\rceil$ \Comment{$blocks=0$ when $q=0$}
			\State $extra \gets \max(0, blocks - 1)$
			\State $best \gets \max(best,\ dp[\ell][q] - extra \cdot c)$
			\EndIf
			\EndFor
			\EndFor
			\State \Return $best$
		\end{algorithmic}
	\end{algorithm}

	\subsection{6. Time complexity}
	Let $Q=\sum_{i=1}^n k_i$ (so $Q \le nK$). The DP fills $n$ stages, each iterating over all $(\ell,q)$ pairs.
	Thus time complexity is $O(n\cdot L \cdot Q)$ and space complexity is $O(L\cdot Q)$.

	Using $Q \le nK$, a valid upper bound is $O(n^2 L K)$. \ \ :contentReference[oaicite:5]{index=5}

\end{document}
